% ============================================================================
%  CAPITULO IV: CONCLUSIONES
% ============================================================================

\section{Conclusiones}

\subsection{Resultado del DNS}

La configuraci\'on del DNS del dominio mediante el Controlador de Dominio de
\textbf{Samba 4} (backend \textbf{SAMBA\_INTERNAL}) permiti\'o resolver de
manera estable la zona autoritativa \texttt{sudoers.lan}. Las consultas desde
el servidor y desde un cliente de la red devolvieron los registros A de los
servicios, los registros SRV de descubrimiento (LDAP, Kerberos) y las
resoluciones externas a trav\'es del \textit{forwarder}, lo que confirma que
la integraci\'on entre el directorio y el DNS funciona correctamente y que los
clientes descubren el dominio sin configuraci\'on manual.

\subsection{Resultado del servidor WEB}

El servidor web (Nginx en contenedor) respondi\'o por el portal y por cada
\textit{virtual host}: \texttt{sudoers.lan} / \texttt{web.sudoers.lan} /
\texttt{www.sudoers.lan} sirvieron el sitio por \textbf{HTTPS} con c\'odigo
\texttt{200 OK}, y los \textit{proxies} hacia impresi\'on, Portainer y las
m\'aquinas de los administradores enrutaron correctamente por nombre. La
redirecci\'on autom\'atica de HTTP a HTTPS y el certificado \textit{wildcard}
de la CA interna verificaron la integraci\'on entre el DNS y el sitio
publicado.

\subsection{Servicios de los administradores}

La publicaci\'on de los subdominios \texttt{david.sudoers.lan} y
\texttt{nicolas.sudoers.lan} demostr\'o la arquitectura del proxy inverso: el
DNS resuelve siempre al punto \'unico de entrada (\texttt{192.168.0.2}) y el
proxy enruta hacia los backends. De esta forma, los servicios de las m\'aquinas
de los admins (\texttt{.51}/\texttt{.52}:8080) quedan expuestos con un nombre
del dominio y el puerto est\'andar 443, sin abrir puertos extra ni depender de
la direcci\'on de cada m\'aquina.

\subsection{Decisiones de dise\~no validadas}

La tabla \ref{tab:decisiones} verific\'o tres principios clave del proyecto:

\begin{itemize}
  \item \textbf{El DC nativo es el \'unico DNS del dominio:} no compite con
        BIND ni con otro servidor; el AD y su zona conviven en el mismo
        demonio \citep{sambateam2024}.
  \item \textbf{Una \'unica puerta web:} todo nombre resuelve al proxy y \'el
        decide el backend; agregar un servicio nuevo no exige cambiar nombres
        ni clientes.
  \item \textbf{HTTPS desde el inicio:} la CA interna cifra toda la intranet
        sin depender de certificados externos; los clientes solo confian en
        \texttt{ca.crt}.
\end{itemize}

\subsection{S\'intesis}

Qued\'o verificado el ciclo completo \textbf{DNS + WEB} en el entorno de
laboratorio: zona del dominio autoritativa, registros de descubrimiento,
resoluci\'on externa, portal publicado por HTTPS y servicios de los miembros
expuestos por subdominio. El trabajo deja las bases para fases posteriores:
activar el webmail y el monitoreo sobre el mismo proxy, incorporar DNSSEC y
evaluar una CA p\'ublica o la renovaci\'on autom\'atica de certificados.